\subsubsection{The finite-projective-dimension triangular row cases}

\begin{lemma}\label{B:six-pd2-row-pairs}
Let $A$ be either the five-dimensional Nakayama algebra of
Lemma~\ref{B:five-exceptional-targets} times $\mathbb Q$, or the
six-dimensional algebra with vertex idempotents $e,f$, radical
basis $x,a,b,c$, and
\[
 x\in eAe,\quad a\in eAf,\quad b\in fAe,\quad c\in eAe,
 \qquad x^2=ab=c,
\]
all other products of radical basis elements being zero.
Let $V$ be the two-dimensional left module on which $b$ is the
nonzero cross-arrow and all other radical elements act by zero.
In the direct-product case the isolated rational factor acts by zero.
Suppose a five-dimensional right module $W$ satisfies
\[
 (\dim\operatorname{Tor}_0^A(W,V),
   \dim\operatorname{Tor}_1^A(W,V),
   \dim\operatorname{Tor}_2^A(W,V))=(2,0,1).
\]
Then $W$ and $V$ have the same annihilator, and are faithful over
the quotient by that annihilator, which is the path algebra of one
arrow.  In the Nakayama product case the annihilator is
$(a,c)\oplus\mathbb Q$; in the six-dimensional case it is $(x,a,c)$.
\end{lemma}
\begin{proof}
In the Nakayama product case write $W=W_N\oplus\mathbb Q^r$.
The Tor table in Lemma~\ref{B:five-exceptional-targets} forces
\[
 W_N=T_2\oplus(e_1A)^{\oplus p}\oplus(e_2A)^{\oplus q}.
\]
The tensor and total dimensions give $p+q=2$ and
$1+3p+2q+r=5$, hence $p=r=0$ and $q=2$.
The faithfulness assertion is the same explicit one-arrow action
as in Lemma~\ref{B:five-two-by-five-pairs}.

For the six-dimensional algebra the following is an exact left
projective resolution, with maps specified by the images of their
vertex generators:
\[
 0\longrightarrow Af\xrightarrow{\ b\ }Ae
   \xrightarrow{\ x\ }Ae\longrightarrow V\longrightarrow0.
\]
Indeed, the first map has image $\operatorname{span}(b,c)$,
which is the kernel of right multiplication by $x$; the latter
has image $\operatorname{span}(x,c)$, the kernel of the quotient
to $V$.  Put $u=\dim We$, $v=\dim Wf$, and write $r,s$ for
the ranks of right multiplication by $x,b$, respectively.
After tensoring the resolution, the three homology dimensions are
\[
 u-r=2,\qquad u-r-s=0,\qquad v-s=1,
 \qquad u+v=5.
\]
Thus $(u,v,r,s)=(2,3,0,2)$.  In particular $x$ acts by zero and
$b:Wf\to We$ is surjective.  The relation $ba=0$ forces $a$ to
act by zero as well, and $c=x^2$ acts by zero.
The surviving $b$-map has rank two; its two vertex projections
and its cross-arrow are linearly independent operators.
Consequently $W$ is faithful over $A/(x,a,c)$, as is $V$.
\end{proof}

\begin{proposition}\label{B:six-pd2-actual-bases}
None of the $124$ finite-projective-dimension triangular row
modules retained by the six-dimensional quotient census can
occur in a monoid with a $3\times6$ rectangular-band minimal
ideal and classifying space $S^2\times S^2$.
\end{proposition}
\begin{proof}
The complete row-module reduction is recorded below in
Proposition~\ref{B:six-row-census}.  Its $124$ cases comprise
$104$ modules over a Nakayama algebra times $\mathbb Q$, and
$20$ modules over the six-dimensional algebra in
Lemma~\ref{B:six-pd2-row-pairs}.  Their row modules are precisely
the modules $V$ in that lemma.  These identifications are checked
by an explicit rational change of algebra basis, followed by all
$36$ multiplication equations, rather than by quiver shape alone.

Here is the additional check on the distinguished monoid bases.
For each basis, form the subspace $L$ spanned by basis elements
whose $V$-operator is zero and by differences of basis elements
having the same $V$-operator.  Direct multiplication verifies
that $L$ is an ideal.  In $98$ Nakayama product cases and in all
$20$ six-dimensional cases, $L$ is the full common annihilator.
In the remaining six Nakayama product cases, $L=(a,c)$; the
isolated rational factor survives.  Thus in every case $A/L$ is
either the path algebra of one arrow or its product with
$\mathbb Q$.

The exact data are in
\path{code/outputs/b_six_pd2_annihilators.json}.
The independent checker
\path{code/b_verify_six_pd2_bases.py} reconstructs all the stated
algebra bases, verifies their complete multiplication, verifies
the row-module actions, and checks each spanning relation and
ideal inclusion.  Its results, including every rational basis
change, are in
\path{code/outputs/b_six_pd2_bases_verified.json}.
It certifies $98$, $20$, and $6$ cases in the three groups above.

To pass from algebra to monoids, let $H$ be the faithful
transformation-pair image of the original monoid.  By the preceding
lemma $W$ and $V$ are faithful over the same algebra quotient.
Thus two distinguished basis elements equal on $V$ are equal on
$W$, and their common augmentation operators are either both
nonzero or both zero.  Equal nonzero augmentation operators imply
equal set maps; zero operators on both sides give constant maps,
whose profile lies in the minimal rectangular band.  It follows
that the contracted algebra of $H$ kills every generator of $L$.
There is no mixed zero/nonzero case here.

Every quotient of the particular algebra $A/L$ is hereditary:
there is at most one arrow, with distinct endpoints, and an
optional isolated vertex.  The band--Tor formula therefore gives
$H_4(BH;\mathbb Q)=0$.  The row and column actions give the same
balanced tensor presentation before and after taking the faithful
image, so $H_2(BM;\mathbb Q)\to H_2(BH;\mathbb Q)$ is an
isomorphism.  All degree-two classes pull back from $BH$, and
all their products vanish.  This excludes $S^2\times S^2$.
\end{proof}
